\section{Core-Value Hierarchy Construction}\label{sec:buildhier}

The core numbers define a nested family of vertex sets.
%
For any level \(k\), the super-level set is \(V_k=\{v\mid \kappa_s(v)\ge k\}\).
%
As \(k\) decreases, vertices enter the set.
%
If connectivity is measured by input-graph edges inside \(G[V_k]\), components are born and merge.
%
\buildHier{} constructs the elder-rule join tree of this edge-connected core-value filtration.
%
This is a useful hierarchy of the computed core values, but it is not claimed to be the exact \(s\)-clique-connected nucleus hierarchy from Definition~\ref{def:nucleus}.

\peelH{} already computes the order needed for this filtration.
%
It removes vertices in nondecreasing \(\kappa_s\).
%
Scanning that deletion log backward inserts vertices in nonincreasing \(\kappa_s\).
%
A union-find structure can then maintain the edge-connected components of the current super-level set.

\stitle{Setup.}
Let \(L=(v_1,k_1),\ldots,(v_n,k_n)\) be the deletion log produced by \peelH{}.
%
Here \(k_i=\kappa_s(v_i)\) and \(k_1\le k_2\le\cdots\le k_n\).
%
No explicit sort is needed.

\stitle{Algorithm state.}
\buildHier{} takes the input graph and the sorted deletion log.
%
Its mutable state is a union-find structure, a birth level for each component root, and the output tree.
%
Its invariant is that after scanning suffix \(i+1,\ldots,n\), union-find represents the connected components of \(G[\{v_j\mid j>i\}]\).
%
The transition inserts \(v_i\), unions it with already-inserted neighbors whose core value is at least \(k_i\), and applies the elder rule when two components merge.
%
The output is the elder-rule join tree of the edge-connected super-level filtration.

\begin{algorithm}[t]
\algfontsize
\caption{\buildHier}
\label{alg:buildhier}
\KwIn{Graph \(G\), deletion log \(L=(v_1,k_1),\ldots,(v_n,k_n)\) with \(k_i\) nondecreasing.}
\KwOut{Elder-rule join tree \(T\) of the edge-connected super-level filtration \(G[V_k]\).}

initialise union-find with each \(v_i\) as a singleton;\quad \(\text{birth}[v_i]\gets k_i\)\;
\(T\gets\) empty tree with one leaf per \(v_i\)\;
\For{\(i\gets n\) \KwTo \(1\)}{
    \(v\gets v_i\);\quad \(k\gets k_i\)\;
    mark \(v\) as inserted\;
    \ForEach{\(u\in N_G(v)\) with \(\kappa_s(u)\ge k\) and \(u\) already inserted}{
        \(r_v\gets\textsc{Find}(v)\);\quad \(r_u\gets\textsc{Find}(u)\)\;
        \If{\(r_v\ne r_u\)}{
            let \(r_{\mathrm{old}}\) and \(r_{\mathrm{young}}\) be the roots with smaller and larger \(\text{birth}\)\tcp*{elder rule}
            attach \(r_{\mathrm{young}}\) as a child of \(r_{\mathrm{old}}\) in \(T\) at level \(k\)\;
            \(\textsc{Union}(r_v,r_u)\)\;
        }
    }
}
\Return{\(T\)}
\end{algorithm}

\stitle{Correctness.}
The reverse scan processes the same vertex-insertion events as the super-level filtration.
%
When a vertex is inserted, all previously inserted vertices have core value at least the current level.
%
Thus every union corresponds to an edge inside the current super-level graph.
%
Conversely, every edge of \(G[V_k]\) is examined when the later endpoint in the reverse scan is inserted, so every edge-connected merge at level \(k\) is performed.
%
When two components merge, the elder rule preserves the component with earlier birth and attaches the younger component in the join tree~\cite{morozov2007persistence}.
%
Because the log is sorted, no later scan step can change a component's birth level.

\stitle{Boundary.}
The algorithm above uses ordinary graph adjacency.
%
Definition~\ref{def:nucleus} uses \(s\)-clique witness connectivity.
%
These notions coincide for \(s=2\), but need not coincide for \(s>2\).
%
An exact \(s\)-clique-connected hierarchy must union vertices only through surviving \(s\)-clique witnesses, which requires reading the \cpi{} or another compact witness representation.
%
Section~\ref{sec:cs-clique} evaluates this clique-coherent variant as an application study.

\stitle{Complexity.}
The scan touches every vertex once.
%
It examines every input edge at most twice, once from each endpoint.
%
Each edge visit performs a constant number of union-find operations.
%
The total time is \(O((n+m)\alpha(n))\), and the \cpi{} is not read.
%
The bound applies to the edge-connected core-value hierarchy, not to the clique-coherent variant.

\stitle{Consequence.}
The edge-connected core-value hierarchy is not a separate decomposition problem in this pipeline.
%
It is a byproduct of the deletion order.
%
Methods that produce unsorted core values need a separate sort before the same reverse scan~\cite{nucleusSariyuce14}.
